% !TEX program = xelatex
\documentclass[11pt,a4paper]{article}
\usepackage[margin=2.2cm,top=2cm]{geometry}
\usepackage{fontspec}
\usepackage{polyglossia}
\setdefaultlanguage{turkish}
\usepackage{amsmath,amssymb}
\usepackage{graphicx}
\usepackage{booktabs}
\usepackage{enumitem}
\usepackage{multicol}
\usepackage[most]{tcolorbox}
\usepackage{fancyhdr}
\usepackage{titlesec}
\usepackage{xcolor}
\usepackage{tikz}
\usepackage{needspace}
\usepackage[hidelinks]{hyperref}

\definecolor{anamavi}{RGB}{0,45,98}
\definecolor{cozumyesil}{RGB}{0,110,50}
\definecolor{uyarikirmizi}{RGB}{180,20,20}

\titleformat{\section}{\large\bfseries\color{anamavi}}{\thesection.}{0.5em}{}
\titleformat{\subsection}{\normalsize\bfseries\color{anamavi}}{\thesubsection}{0.5em}{}

\pagestyle{fancy}\fancyhf{}
\lhead{\small MMAT105 Genel Matematik}
\rhead{\small Hafta 6 -- Ders Notu}
\cfoot{\small\thepage}
\renewcommand{\headrulewidth}{0.4pt}

\newtcolorbox{kural}[1][Kural]{colback=anamavi!5,colframe=anamavi,title=#1,fonttitle=\bfseries,breakable,left=4pt,right=4pt,top=2pt,bottom=2pt}
\newtcolorbox{dikkat}[1][Dikkat]{colback=uyarikirmizi!5,colframe=uyarikirmizi,title=#1,fonttitle=\bfseries,breakable,left=4pt,right=4pt,top=2pt,bottom=2pt}
\newcounter{ornek}
\newenvironment{ornek}{\refstepcounter{ornek}\par\medskip\noindent\textbf{\color{anamavi}Örnek \theornek.}\ }{\par}
\newenvironment{cozum}{\begin{tcolorbox}[colback=cozumyesil!6,colframe=cozumyesil,boxrule=0.5pt,left=4pt,right=4pt,top=2pt,bottom=2pt,breakable]\textbf{\color{cozumyesil}Çözüm.}\ }{\end{tcolorbox}}
\newcounter{alis}
\newenvironment{alistirma}{\refstepcounter{alis}\par\medskip\needspace{9\baselineskip}\noindent\textbf{\color{anamavi}Alıştırma \thealis.}\ }{\par}
\newcommand{\kendinizoncesi}{\needspace{12\baselineskip}}
\newcommand{\ab}[1]{\overline{#1}}
\setlength{\parindent}{0pt}
\setlength{\parskip}{3pt}

\begin{document}

% ---------------- Başlık ----------------
\begin{minipage}[c]{\textwidth}
{\LARGE\bfseries\color{anamavi} MMAT105 Genel Matematik}\\[4pt]
{\large Hafta 6 Ders Notu: Birinci Dereceden\\ İki Bilinmeyenli Denklemler ve Denklem Sistemleri}\\[6pt]
{\small 2026--2027 Güz}
\end{minipage}

\vspace{4pt}\noindent\textcolor{anamavi}{\rule{\textwidth}{1pt}}

\textbf{Bu notun kullanımı.} Her bölümde önce kısa bir kural kutusu, ardından çözümlü örnekler yer alır. Örnekleri çözümüne bakmadan önce kendiniz deneyin. Notun sonunda önce çözümlü alıştırmalar, ardından cevap anahtarlı ``Kendiniz Çözün'' alıştırmaları yer alır.

\textbf{Altın kural.} Bulduğunuz \((x,y)\) ikilisini her zaman \textbf{iki denklemde de} kontrol edin.

% =====================================================
\section{İki Bilinmeyenli Denklem}

\begin{kural}[Tanım]
\(a,b,c\) gerçek sayılar, \(a\neq 0\) ve \(b\neq 0\) olmak üzere \(ax+by+c=0\) biçimindeki denkleme \textbf{birinci dereceden iki bilinmeyenli denklem} denir.
\begin{itemize}[leftmargin=*,itemsep=1pt]
\item Denklemi sağlayan her \((x,y)\) sıralı ikilisi bir \textbf{çözümdür}. \((x,y)\) sıralıdır: \((2,5)\neq(5,2)\).
\item Tek bir denklemin sonsuz sayıda çözümü vardır: \(x\)'e istenen değer verilir, \(y\) hesaplanır.
\item Bir ikilinin çözüm olup olmadığı, değerler yerine konup eşitliğin sağlanıp sağlanmadığına bakılarak anlaşılır.
\end{itemize}
\end{kural}

\begin{ornek}
\(3x-2y=6\) denkleminde \(x=4\) ise \(y\) kaçtır? \(y=0\) ise \(x\) kaçtır?
\end{ornek}
\begin{cozum}
\(x=4\): \ \(12-2y=6\Rightarrow 2y=6\Rightarrow y=3\). \qquad \(y=0\): \ \(3x=6\Rightarrow x=2\).\\
Böylece \((4,3)\) ve \((2,0)\) bu denklemin iki çözümüdür.
\end{cozum}

\begin{ornek}
\((m+1,\ 2m)\) ikilisi \(x+2y=11\) denklemini sağlıyorsa \(m\) kaçtır?
\end{ornek}
\begin{cozum}
\[(m+1)+2\cdot 2m=11\ \Rightarrow\ 5m+1=11\ \Rightarrow\ m=2\]
İkili: \((3,4)\). Kontrol: \(3+8=11\).
\end{cozum}

\begin{ornek}
\(3x+5y=30\) denklemini sağlayan kaç tane \((x,y)\) doğal sayı ikilisi vardır? \ (\(\mathbb{N}=\{0,1,2,\dots\}\))
\end{ornek}
\begin{cozum}
\(3x=30-5y\) ifadesi \(3\)'e bölünmeli; \(30\) \(3\)'e bölündüğünden \(5y\), dolayısıyla \(y\) \(3\)'ün katı olmalı. \(5y\le 30\Rightarrow y\le 6\).\\
\(y=0,3,6\ \Rightarrow\ x=10,5,0\). \quad İkililer: \((10,0),(5,3),(0,6)\) \(\Rightarrow\) \(3\) tane.
\end{cozum}

% =====================================================
\section{Denklem Sistemi ve Yok Etme Yöntemi}

\begin{kural}[Yok etme yöntemi]
\(\begin{cases}a_1x+b_1y=c_1\\ a_2x+b_2y=c_2\end{cases}\) sisteminin çözümü, \textbf{iki denklemi birlikte} sağlayan \((x,y)\) ikilisidir.
\begin{enumerate}[leftmargin=*,itemsep=1pt]
\item Bir bilinmeyenin katsayılarını, denklemleri uygun sayılarla çarparak eşit ya da zıt yap.
\item Katsayılar zıt ise taraf tarafa \textbf{topla}, eşit ise taraf tarafa \textbf{çıkar}: bir bilinmeyen yok olur.
\item Bulunan değeri denklemlerden birinde yerine koy; ikiliyi iki denklemde de kontrol et.
\end{enumerate}
\end{kural}

\begin{dikkat}[Sık yapılan hatalar]
Denklemi çarparken sağ tarafı da çarpın. Taraf tarafa çıkarırken ikinci denklemin \textbf{her terimi} işaret değiştirir: \((2x+3y)-(2x-y)=4y\).
\end{dikkat}

\begin{ornek}
\(\begin{cases}2x+y=10\\ x-y=2\end{cases}\) sisteminin çözümü nedir?
\end{ornek}
\begin{cozum}
\(y\)'lerin katsayıları zıt; topla: \(3x=12\Rightarrow x=4\). \quad \(4-y=2\Rightarrow y=2\). \quad Çözüm: \((4,2)\).\\
Kontrol: \(8+2=10\) \checkmark, \ \(4-2=2\) \checkmark
\end{cozum}

\begin{ornek}
\(\begin{cases}3x-5y=-1\\ 2x+3y=12\end{cases}\) sisteminin çözümü nedir?
\end{ornek}
\begin{cozum}
\(x\)'i yok edelim: 1. denklem \(\times 2\), 2. denklem \(\times 3\):
\[6x-10y=-2,\qquad 6x+9y=36\]
Çıkar (2. \(-\) 1.): \(19y=38\Rightarrow y=2\). \quad \(2x+6=12\Rightarrow x=3\). \quad Çözüm: \((3,2)\).
\end{cozum}

\begin{ornek}
\(\begin{cases}4x+3y=2\\ 5x-2y=14\end{cases}\) ise \(x-y\) kaçtır?
\end{ornek}
\begin{cozum}
\(y\)'yi yok edelim: 1. denklem \(\times 2\), 2. denklem \(\times 3\):
\[8x+6y=4,\qquad 15x-6y=42\ \Rightarrow\ 23x=46\ \Rightarrow\ x=2\]
\(8+3y=2\Rightarrow y=-2\). \quad \(x-y=2-(-2)=4\).
\end{cozum}

% =====================================================
\section{Yerine Koyma Yöntemi; Kesirli ve Ondalıklı Katsayılar}

\begin{kural}[Yerine koyma]
Bir denklemde bir bilinmeyeni yalnız bırak (\(y=\dots\) ya da \(x=\dots\)), bu ifadeyi \textbf{parantez içinde} diğer denklemde yerine koy, tek bilinmeyenli denklemi çöz, sonra geri dön.\\
\textbf{Kesirli katsayılar:} Her denklemi kendi paydalarının EKOK'u ile çarp. \ \textbf{Ondalıklı katsayılar:} \(10\), \(100\), \dots\ ile çarp.
\end{kural}

\needspace{8\baselineskip}
\begin{ornek}
\(\begin{cases}2x+y=7\\ 3x-2y=0\end{cases}\) sistemini yerine koyma yöntemiyle çözünüz.
\end{ornek}
\begin{cozum}
1. denklemden \(y=7-2x\). \ 2. denklemde yerine koyalım:
\[3x-2(7-2x)=0\ \Rightarrow\ 3x-14+4x=0\ \Rightarrow\ 7x=14\ \Rightarrow\ x=2\]
\(y=7-4=3\). \quad Çözüm: \((2,3)\).
\end{cozum}

\begin{ornek}
\(\dfrac{x+1}{3}=\dfrac{y-1}{2}\) \ ve \ \(x+y=15\) olduğuna göre \(x\) ve \(y\) kaçtır?
\end{ornek}
\begin{cozum}
İçler--dışlar çarpımı: \(2(x+1)=3(y-1)\Rightarrow 2x-3y=-5\). \ İkinci denklemden \(x=15-y\):
\[2(15-y)-3y=-5\ \Rightarrow\ 30-5y=-5\ \Rightarrow\ y=7,\quad x=8\]
Kontrol: \(\dfrac{9}{3}=3=\dfrac{6}{2}\) \checkmark
\end{cozum}

\begin{ornek}
\(\begin{cases}0{,}5x+0{,}2y=2{,}6\\ x-0{,}4y=2{,}8\end{cases}\) sisteminin çözümü nedir?
\end{ornek}
\begin{cozum}
1. denklem \(\times 10\): \(5x+2y=26\). \ 2. denklem \(\times 5\): \(5x-2y=14\).
\[\text{Topla: } 10x=40\Rightarrow x=4,\qquad 20+2y=26\Rightarrow y=3\]
Çözüm: \((4,3)\).
\end{cozum}

% =====================================================
\section{Kısa Yollar: Simetrik Sistemler ve İstenen İfadeler}

\begin{kural}[Simetrik sistem]
\(\begin{cases}ax+by=m\\ bx+ay=n\end{cases}\) \((a\neq b,\ a+b\neq 0)\) biçimindeki sistemlerde:
\[\text{topla: } x+y=\frac{m+n}{a+b},\qquad \text{çıkar: } x-y=\frac{m-n}{a-b}\]
Soru \(x+y\), \(x-y\) veya \(x\cdot y\) istiyorsa önce denklemleri toplamayı ya da çıkarmayı deneyin.
\end{kural}

\begin{ornek}
\(\begin{cases}6x+5y=39\\ 5x+6y=38\end{cases}\) olduğuna göre \(x\cdot y\) kaçtır?
\end{ornek}
\begin{cozum}
Topla: \(11(x+y)=77\Rightarrow x+y=7\). \quad Çıkar: \(x-y=1\).\\
\(2x=8\Rightarrow x=4\), \ \(y=3\). \quad \(x\cdot y=12\).
\end{cozum}

\begin{ornek}
\textbf{(İleri örnek: üç bilinmeyen.)} \(x+y=7\), \ \(y+z=9\), \ \(x+z=10\) ise \(x\), \(y\), \(z\) kaçtır?
\end{ornek}
\begin{cozum}
Üç denklemi toplayalım: \(2(x+y+z)=26\Rightarrow x+y+z=13\).\\
\(z=13-(x+y)=13-7=6\), \quad \(x=13-9=4\), \quad \(y=13-10=3\).
\end{cozum}

% =====================================================
\section{Özel Durumlar ve Grafik Yorumu}

\begin{kural}[Çözüm sayısı]
\(a_1x+b_1y=c_1\) ve \(a_2x+b_2y=c_2\) sisteminde (oranlar tanımlı olmak üzere):
\begin{center}
\(\renewcommand{\arraystretch}{1.2}
\begin{array}{lll}
\text{Koşul} & \text{Çözüm} & \text{Grafik}\\\hline
\dfrac{a_1}{a_2}\neq\dfrac{b_1}{b_2} & \text{tek çözüm} & \text{kesişen iki doğru}\\[10pt]
\dfrac{a_1}{a_2}=\dfrac{b_1}{b_2}\neq\dfrac{c_1}{c_2} & \text{çözüm yok} & \text{paralel iki doğru}\\[10pt]
\dfrac{a_1}{a_2}=\dfrac{b_1}{b_2}=\dfrac{c_1}{c_2} & \text{sonsuz çözüm} & \text{çakışık (aynı) doğru}
\end{array}\)
\end{center}
Paydalardan biri \(0\) ise oran yerine içler--dışlar çarpımı kullanılır: tek çözüm \(\iff a_1b_2-a_2b_1\neq 0\).
\end{kural}

\begin{center}
\begin{tikzpicture}[scale=0.5]
\begin{scope}
\draw[->,gray] (-0.5,0)--(5,0) node[below]{\scriptsize \(x\)};
\draw[->,gray] (0,-1.5)--(0,5) node[left]{\scriptsize \(y\)};
\draw[thick,anamavi] (-0.5,4.5)--(4.5,-0.5);
\draw[thick,uyarikirmizi] (1,-1)--(5,3);
\fill (3,1) circle (3pt); \node[above right] at (3,1) {\scriptsize \((3,1)\)};
\node[align=center] at (2.3,-2.8) {\small kesişen (tek çözüm)\\ \small \(x+y=4,\ x-y=2\)};
\end{scope}
\begin{scope}[xshift=10cm]
\draw[->,gray] (-0.5,0)--(5,0) node[below]{\scriptsize \(x\)};
\draw[->,gray] (0,-1.5)--(0,5) node[left]{\scriptsize \(y\)};
\draw[thick,anamavi] (-0.5,4.5)--(4.5,-0.5);
\draw[thick,uyarikirmizi] (-0.5,1.5)--(2.5,-1.5);
\node[align=center] at (2.3,-2.8) {\small paralel (çözüm yok)\\ \small \(x+y=4,\ x+y=1\)};
\end{scope}
\begin{scope}[xshift=20cm]
\draw[->,gray] (-0.5,0)--(5,0) node[below]{\scriptsize \(x\)};
\draw[->,gray] (0,-1.5)--(0,5) node[left]{\scriptsize \(y\)};
\draw[line width=3pt,uyarikirmizi!50] (-0.5,4.5)--(4.5,-0.5);
\draw[thick,anamavi] (-0.5,4.5)--(4.5,-0.5);
\node[align=center] at (2.3,-2.8) {\small çakışık (sonsuz çözüm)\\ \small \(x+y=4,\ 2x+2y=8\)};
\end{scope}
\end{tikzpicture}
\end{center}

\begin{ornek}
\(\begin{cases}3x-y=2\\ 9x+ay=6\end{cases}\) sisteminin sonsuz çözümü varsa \(a\) kaçtır?
\end{ornek}
\begin{cozum}
\(\dfrac{3}{9}=\dfrac{-1}{a}=\dfrac{2}{6}=\dfrac13\ \Rightarrow\ a=-3\). \ (2. denklem, 1. denklemin \(3\) katıdır.)
\end{cozum}

\begin{ornek}
\(\begin{cases}ax+2y=5\\ 3x-y=4\end{cases}\) sisteminin çözümü yoksa \(a\) kaçtır?
\end{ornek}
\begin{cozum}
\(\dfrac{a}{3}=\dfrac{2}{-1}\Rightarrow a=-6\). \ Kontrol: \(\dfrac{c_1}{c_2}=\dfrac54\neq -2\) \checkmark \ (\(a=-6\) için \(-6x+2y=5\) ile \(-6x+2y=-8\) çelişir.)
\end{cozum}

\begin{ornek}
\(\begin{cases}2x+3y=5\\ 4x+6y=k\end{cases}\) sistemi için: sonsuz çözüm varsa \(k\) kaçtır? Çözüm yoksa \(k\) için ne söylenebilir?
\end{ornek}
\begin{cozum}
\(\dfrac24=\dfrac36=\dfrac12\) olduğundan sistem hiçbir \(k\) için tek çözümlü değildir.\\
Sonsuz çözüm: \(\dfrac{5}{k}=\dfrac12\Rightarrow k=10\). \quad Çözüm yok: \(k\neq 10\).
\end{cozum}

% =====================================================
\section{Sözel Problemler}

\begin{kural}[Sistem kurma]
\begin{enumerate}[leftmargin=*,itemsep=1pt]
\item İki bilinmeyene harf verin ve ne olduklarını yazın.
\item Metindeki iki bağımsız bilgiyi iki denkleme çevirin (ör. adet: \(x+y=\dots\); \ tutar: \(px+qy=\dots\)).
\item Sistemi çözün; \textbf{sorulan} büyüklüğü verin ve sonucu metinde kontrol edin.
\end{enumerate}
İki basamaklı sayılar: \(\ab{ab}=10a+b\) \((a\neq 0)\), \ \(\ab{ab}+\ab{ba}=11(a+b)\), \ \(\ab{ab}-\ab{ba}=9(a-b)\).
\end{kural}

\begin{ornek}
İki sayının farkı \(12\)'dir. Büyük sayının yarısı, küçük sayıdan \(1\) fazladır. Bu sayılar nedir?
\end{ornek}
\begin{cozum}
Büyük \(x\), küçük \(y\): \ \(x-y=12\), \ \(\dfrac{x}{2}=y+1\Rightarrow x=2y+2\).
\[(2y+2)-y=12\Rightarrow y=10,\quad x=22\qquad(\text{kontrol: } 11=10+1)\]
\end{cozum}

\begin{ornek}
Bir kumbarada yalnızca \(1\) TL'lik ve \(50\) kuruşluk madeni paralardan toplam \(40\) tane ve \(29\) TL vardır. Kaç tane \(50\) kuruşluk vardır?
\end{ornek}
\begin{cozum}
\(x\): 1 TL'lik, \(y\): 50 kuruşluk. \ \(x+y=40\), \ \(x+0{,}5y=29\).\\
Çıkar: \(0{,}5y=11\Rightarrow y=22\), \ \(x=18\). \quad Kontrol: \(18+11=29\) \checkmark
\end{cozum}

\begin{ornek}
Bir sinemada yetişkin bileti \(120\) TL, çocuk bileti \(80\) TL'dir. Bir grup \(14\) bilet için \(1440\) TL ödemiştir. Grupta kaç yetişkin vardır?
\end{ornek}
\begin{cozum}
\(x+y=14\), \ \(120x+80y=1440\Rightarrow 3x+2y=36\). \ \(3x+2(14-x)=36\Rightarrow x=8\), \ \(y=6\).\\
Kontrol: \(960+480=1440\) \checkmark \quad Cevap: \(8\) yetişkin.
\end{cozum}

\begin{ornek}
\(\ab{ab}\) iki basamaklı bir sayıdır. \(\ab{ab}+\ab{ba}=110\) ve \(a-b=4\) ise \(\ab{ab}\) kaçtır?
\end{ornek}
\begin{cozum}
\(11(a+b)=110\Rightarrow a+b=10\); \ \(a-b=4\). \ Topla: \(2a=14\Rightarrow a=7\), \ \(b=3\). \quad \(\ab{ab}=73\) \ (kontrol: \(73+37=110\)).
\end{cozum}

\begin{ornek}
Bir çiftlikteki ördek ve ineklerin toplam \(35\) başı ve \(94\) ayağı vardır. Kaç inek vardır?
\end{ornek}
\begin{cozum}
\(o+i=35\), \ \(2o+4i=94\Rightarrow o+2i=47\). \ Çıkar: \(i=12\), \ \(o=23\).\\
Kontrol: \(46+48=94\) \checkmark
\end{cozum}

\begin{ornek}
Bir anne ile kızının bugünkü yaşları toplamı \(48\)'dir. Annenin yaşı kızının yaşının \(3\) katıdır. Yaşları arasındaki fark kaçtır?
\end{ornek}
\begin{cozum}
\(a+k=48\), \ \(a=3k\Rightarrow 4k=48\Rightarrow k=12\), \ \(a=36\). \quad Fark: \(24\). \ (Yaş problemleri 13. haftada ayrıntılı işlenecektir.)
\end{cozum}

% =====================================================
\section{Özet: Formül Kartı}
\begin{center}
\renewcommand{\arraystretch}{1.45}
\begin{tabular}{@{}lp{0.68\textwidth}@{}}
\toprule
Konu & Formül / Kural\\\midrule
İki bilinmeyenli denklem & \(ax+by+c=0\) \((a,b\neq 0)\); çözüm sıralı ikili \((x,y)\); sonsuz çözüm\\
Yok etme & katsayıları eşitle/zıtla; eşitse çıkar, zıtsa topla\\
Yerine koyma & \(y=\dots\) yaz, diğer denklemde \textbf{parantezle} yerine koy\\
Kesir / ondalık & paydaların EKOK'u ile ya da \(10,100,\dots\) ile çarp\\
Simetrik sistem & topla \(\Rightarrow x+y\), \ çıkar \(\Rightarrow x-y\)\\
Tek çözüm & \(\dfrac{a_1}{a_2}\neq\dfrac{b_1}{b_2}\) \ (kesişen doğrular)\\
Çözüm yok & \(\dfrac{a_1}{a_2}=\dfrac{b_1}{b_2}\neq\dfrac{c_1}{c_2}\) \ (paralel doğrular)\\
Sonsuz çözüm & \(\dfrac{a_1}{a_2}=\dfrac{b_1}{b_2}=\dfrac{c_1}{c_2}\) \ (çakışık doğrular)\\
İki basamaklı sayı & \(\ab{ab}=10a+b\); \ \(\ab{ab}+\ab{ba}=11(a+b)\), \ \(\ab{ab}-\ab{ba}=9(a-b)\)\\
Adet -- tutar & \(x+y=\text{adet}\), \ \(px+qy=\text{toplam tutar}\)\\\bottomrule
\end{tabular}
\end{center}

% =====================================================
\needspace{15\baselineskip}
% Hafta 6 -- alıştırma gövdesi (ders notu ve ayrı alıştırma kağıdı tarafından \input edilir)
\section{Çözümlü Alıştırmalar}
Her soruyu önce kendiniz deneyin, sonra çözümü okuyun.

\begin{alistirma}
\(4x+y=12\) denklemini sağlayan kaç tane \((x,y)\) doğal sayı ikilisi vardır? \ (\(\mathbb{N}=\{0,1,2,\dots\}\))
\end{alistirma}
\begin{cozum}
\(y=12-4x\). \(y\) doğal sayı olmalı: \(12-4x\ge 0\Rightarrow 4x\le 12\Rightarrow x\le 3\).\\
\(x=0,1,2,3\ \Rightarrow\ y=12,8,4,0\). \quad İkililer: \((0,12),(1,8),(2,4),(3,0)\) \(\Rightarrow\) \(4\) tane.
\end{cozum}

\begin{alistirma}
\(\begin{cases}3x+2y=7\\ 5x-3y=-1\end{cases}\) sisteminin çözümü nedir?
\end{alistirma}
\begin{cozum}
\(y\)'yi yok edelim: 1. denklem \(\times 3\), 2. denklem \(\times 2\):
\[9x+6y=21,\qquad 10x-6y=-2\]
Taraf tarafa toplayalım: \(19x=19\Rightarrow x=1\). \quad 1. denklemde: \(3+2y=7\Rightarrow 2y=4\Rightarrow y=2\).\\
Çözüm: \((1,2)\). \quad Kontrol: \(5\cdot 1-3\cdot 2=-1\) \checkmark
\end{cozum}

\begin{alistirma}
\(\begin{cases}\dfrac{x}{3}-\dfrac{y}{2}=0\\[6pt] \dfrac{x}{2}+\dfrac{y}{4}=4\end{cases}\) sisteminin çözümü nedir?
\end{alistirma}
\begin{cozum}
Paydaları temizleyelim: 1. denklem \(\times 6\): \(2x-3y=0\); \quad 2. denklem \(\times 4\): \(2x+y=16\).\\
\(x\)'lerin katsayıları eşit; çıkaralım: \((2x+y)-(2x-3y)=16-0\Rightarrow 4y=16\Rightarrow y=4\).\\
\(2x+4=16\Rightarrow x=6\). \quad Çözüm: \((6,4)\). \quad Kontrol: \(\dfrac63-\dfrac42=0\) \checkmark, \ \(\dfrac62+\dfrac44=4\) \checkmark
\end{cozum}

\begin{alistirma}
\(4x+3y=18\) ve \(3x+4y=17\) olduğuna göre \(x+y\) ve \(x\cdot y\) kaçtır?
\end{alistirma}
\begin{cozum}
Simetrik sistem. Toplayalım: \(7x+7y=35\Rightarrow x+y=5\).\\
Çıkaralım: \(x-y=18-17=1\).\\
\(x+y=5\) ve \(x-y=1\) \(\Rightarrow\) \(2x=6\Rightarrow x=3\), \(y=2\). \quad \(x\cdot y=6\).
\end{cozum}

\begin{alistirma}
\(\begin{cases}2x+ay=3\\ 6x+9y=b\end{cases}\) sisteminin sonsuz çözümü varsa \(a+b\) kaçtır?
\end{alistirma}
\begin{cozum}
Sonsuz çözüm \(\Rightarrow\) \(\dfrac{2}{6}=\dfrac{a}{9}=\dfrac{3}{b}\). \ İlk oran \(\dfrac13\) olduğundan:
\[\frac{a}{9}=\frac13\Rightarrow a=3,\qquad \frac{3}{b}=\frac13\Rightarrow b=9,\qquad a+b=12\]
Kontrol: 1. denklemin \(3\) katı \(6x+9y=9\), yani 2. denklemin kendisi \checkmark
\end{cozum}

\begin{alistirma}
\(\begin{cases}mx-y=4\\ 2x+3y=5\end{cases}\) sisteminin tek çözümü \textbf{yoksa} \(m\) kaçtır? Bu durumda sistemin kaç çözümü vardır?
\end{alistirma}
\begin{cozum}
Tek çözüm yok \(\Rightarrow\) \(\dfrac{m}{2}=\dfrac{-1}{3}\Rightarrow m=-\dfrac23\).\\
Sabitlerin oranı \(\dfrac45\neq-\dfrac13\) olduğundan sistemin çözümü \textbf{yoktur} (\(0\) çözüm).\\
Gerçekten: \(m=-\dfrac23\) için 1. denklem \(\times(-3)\): \(2x+3y=-12\); oysa 2. denklem \(2x+3y=5\): çelişki.
\end{cozum}

\begin{alistirma}
Bir konserde öğrenci bileti \(150\) TL, tam bilet \(250\) TL'dir. Satılan toplam \(40\) biletten \(7400\) TL elde edilmiştir. Kaç öğrenci bileti satılmıştır?
\end{alistirma}
\begin{cozum}
\(x\): öğrenci bileti, \(y\): tam bilet sayısı. \quad \(x+y=40\), \ \(150x+250y=7400\Rightarrow 3x+5y=148\) (\(50\)'ye böldük).\\
\(x=40-y\) yerine koyalım: \(3(40-y)+5y=148\Rightarrow 120+2y=148\Rightarrow y=14\), \ \(x=26\).\\
Kontrol: \(150\cdot 26+250\cdot 14=3900+3500=7400\) \checkmark \quad Cevap: \(26\) öğrenci bileti.
\end{cozum}

\begin{alistirma}
İki basamaklı \(\ab{ab}\) sayısının rakamları toplamı \(12\)'dir ve \(\ab{ab}-\ab{ba}=36\)'dır. \(\ab{ab}\) kaçtır?
\end{alistirma}
\begin{cozum}
\(\ab{ab}-\ab{ba}=(10a+b)-(10b+a)=9(a-b)=36\Rightarrow a-b=4\).\\
\(a+b=12\) ile toplayalım: \(2a=16\Rightarrow a=8\), \ \(b=4\). \quad \(\ab{ab}=84\). \quad Kontrol: \(84-48=36\) \checkmark
\end{cozum}

\kendinizoncesi
\section{Alıştırmalar (Kendiniz Çözün)}

\begin{enumerate}[label=\textbf{\arabic*.},leftmargin=*,itemsep=6pt]
\item \((3,k)\) ikilisi \(2x-y=4\) denklemini sağlıyorsa \(k\) kaçtır?
\item \((-2,1)\), \((2,1)\), \((1,0)\), \((4,-1)\) ikililerinden hangileri \(x+3y=1\) denklemini sağlar?
\item \(x+y=15\) ve \(x-y=7\) ise \(x\) ve \(y\) kaçtır?
\item \(\begin{cases}2x+3y=8\\ 4x-3y=-2\end{cases}\) sisteminin çözümü nedir?
\item \(y=3x-5\) ve \(2x+y=10\) ise \(x\cdot y\) kaçtır?
\item \(\begin{cases}0{,}3x+0{,}2y=1{,}3\\ 0{,}1x-0{,}4y=-0{,}5\end{cases}\) sisteminin çözümü nedir?
\item \(x+y=10\) ve \(x-y=4\) ise \(x\cdot y\) kaçtır?
\item \(\begin{cases}(k+1)x+2y=5\\ 6x+4y=7\end{cases}\) sisteminin çözümü yoksa \(k\) kaçtır?
\item İki sayının toplamı \(63\)'tür. Sayılardan biri diğerinin \(2\) katıdır. Sayıların farkı kaçtır?
\item Bir kasada yalnızca \(10\) TL'lik ve \(20\) TL'lik banknotlardan toplam \(15\) banknot ve \(210\) TL vardır. Kaç tane \(20\) TL'lik banknot vardır?
\item Bir çiftlikteki tavuk ve tavşanların toplam \(26\) başı ve \(76\) bacağı vardır. Kaç tavuk vardır?
\item \(30\) soruluk bir testte doğru cevap \(+4\), yanlış cevap \(-1\) puandır. Bütün soruları cevaplayan bir öğrenci \(80\) puan almıştır. Kaç soruyu doğru cevaplamıştır?
\end{enumerate}

\subsection*{Cevap Anahtarı}
{\small
\begin{multicols}{2}
\begin{enumerate}[label=\textbf{\arabic*.},leftmargin=*,itemsep=1pt]
\item \(6-k=4\Rightarrow k=2\).
\item \((-2,1)\), \((1,0)\), \((4,-1)\) sağlar; \((2,1)\): \(5\neq 1\).
\item Topla: \(2x=22\Rightarrow x=11\), \ \(y=4\).
\item Topla: \(6x=6\Rightarrow x=1\); \ \(y=2\). \ \((1,2)\).
\item \(2x+3x-5=10\Rightarrow x=3\), \(y=4\); \ \(x\cdot y=12\).
\item \(\times 10\): \(3x+2y=13\), \(x=4y-5\) \(\Rightarrow\) \(14y=28\), \(y=2\), \(x=3\). \ \((3,2)\).
\item \(x=7\), \(y=3\) \(\Rightarrow\) \(x\cdot y=21\).
\item \(\dfrac{k+1}{6}=\dfrac24\Rightarrow k=2\); \ \(\dfrac57\neq\dfrac12\) \checkmark
\item \(x=2y\), \(3y=63\Rightarrow y=21\), \(x=42\); fark \(21\).
\item \(x+2y=21\), \(x+y=15\) \(\Rightarrow\) \(y=6\).
\item \(t+2r=38\), \(t+r=26\) \(\Rightarrow\) \(r=12\), \(t=14\).
\item \(d+y=30\), \(4d-y=80\) \(\Rightarrow\) \(5d=110\), \(d=22\).
\end{enumerate}
\end{multicols}}


\vspace{6pt}
\noindent\textcolor{anamavi}{\rule{\textwidth}{0.5pt}}\\
{\small Gelecek hafta: Basit eşitsizlikler ve mutlak değer.}

\end{document}
